\section{From $C_\ell^{AB}(r_1,r_2)$ to $S_\ell^{AB}(k_1,k_2)$ via Hankel transforms}

This note derives the Spherical Fourier-Bessel (SFB) power spectrum in full
generality for two projected fields $f_A$ and $f_B$, and then specializes the
result to galaxy clustering, weak lensing shear, and their cross-correlation.
Every step is written explicitly so that each transformation can be traced.

\subsection{Notation and physical meaning}

We define two observed fields on the light cone:
\begin{equation}
f_X(\hat{\mathbf n},r),\qquad X\in\{A,B\},
\end{equation}
where $r$ is comoving distance and $\hat{\mathbf n}$ is the sky direction.
For each field, the radial dependence enters through a kernel $W^X(r)$.

The symbols used below have the following roles:
\begin{itemize}
\item $\ell$: angular multipole, controlling angular scale on the sky.
\item $j_\ell$: spherical Bessel function, coupling radial and angular modes.
\item $q$: latent 3D Fourier mode over which the theory integrates.
\item $k_1,k_2$: observed SFB radial wavenumbers.
\item $P_{AB}(q)$: 3D power spectrum sourcing correlations between sectors $A,B$.
\item $\alpha_A,\alpha_B$: probe-dependent powers (for example 0 or 2) encoding
the operator structure of each observable.
\item $\mathcal N_\ell^{AB}$: prefactor carrying normalization and spin factors.
\end{itemize}

\subsection{Step 1: Start from the general harmonic-space radial kernel}

We assume the generalized BLAST-like form
\begin{equation}
C_\ell^{AB}(r_1,r_2)
= \mathcal N_\ell^{AB}\,W^A(r_1)W^B(r_2)
\int_0^{\infty} dq\,q^2\,P_{AB}(q)
\frac{j_\ell(qr_1)}{(qr_1)^{\alpha_A}}
\frac{j_\ell(qr_2)}{(qr_2)^{\alpha_B}}.
\label{eq:gen_c_ell_rr}
\end{equation}

What this equation does:
\begin{itemize}
\item It correlates two radial shells $r_1,r_2$ for a fixed angular multipole $\ell$.
\item The radial kernels $W^A,W^B$ encode survey selection and physics weights.
\item The $q$-integral builds the covariance from all 3D modes.
\item The factors $(qr)^{-\alpha_X}$ let one expression cover clustering,
shear, and mixed probes.
\end{itemize}

\subsection{Step 2: Define the double spherical Hankel transform}

To move from radial positions $(r_1,r_2)$ to radial frequencies $(k_1,k_2)$,
define
\begin{equation}
S_\ell^{AB}(k_1,k_2)
\equiv \left(\frac{2}{\pi}\right)^2
\int_0^{\infty} dr_1\,r_1^2
\int_0^{\infty} dr_2\,r_2^2\,
C_\ell^{AB}(r_1,r_2)\,j_\ell(k_1r_1)\,j_\ell(k_2r_2).
\label{eq:def_double_hankel_general}
\end{equation}

What this step does:
\begin{itemize}
\item It is the spherical analogue of a Fourier transform in each radial coordinate.
\item It produces an SFB covariance matrix in $(k_1,k_2)$ at fixed $\ell$.
\item It keeps off-diagonal couplings, so radial mode mixing is preserved.
\end{itemize}

\subsection{Step 3: Substitute Eq.~\eqref{eq:gen_c_ell_rr} into the transform}

Inserting Eq.~\eqref{eq:gen_c_ell_rr} into Eq.~\eqref{eq:def_double_hankel_general},
\begin{align}
S_\ell^{AB}(k_1,k_2)
&= \left(\frac{2}{\pi}\right)^2\mathcal N_\ell^{AB}
\int_0^{\infty} dq\,q^2\,P_{AB}(q)
\Biggl[\int_0^{\infty} dr_1\,r_1^2\,W^A(r_1)
\frac{j_\ell(qr_1)}{(qr_1)^{\alpha_A}}j_\ell(k_1r_1)\Biggr]
\nonumber\\
&\hspace{1.5cm}\times
\Biggl[\int_0^{\infty} dr_2\,r_2^2\,W^B(r_2)
\frac{j_\ell(qr_2)}{(qr_2)^{\alpha_B}}j_\ell(k_2r_2)\Biggr].
\label{eq:inserted_general}
\end{align}

What this step does:
\begin{itemize}
\item It factorizes the expression into one radial response for field $A$ and one
radial response for field $B$.
\item It isolates all cosmology dependence in $P_{AB}(q)$ and the kernels.
\item It prepares the expression for a compact window-function form.
\end{itemize}

\subsection{Step 4: Define generalized radial response kernels}

Define
\begin{equation}
\Psi_{\ell}^{A}(k,q)
\equiv \int_0^{\infty} dr\,r^2\,W^A(r)
\frac{j_\ell(qr)}{(qr)^{\alpha_A}}j_\ell(kr),
\label{eq:psi_A_general}
\end{equation}
\begin{equation}
\Psi_{\ell}^{B}(k,q)
\equiv \int_0^{\infty} dr\,r^2\,W^B(r)
\frac{j_\ell(qr)}{(qr)^{\alpha_B}}j_\ell(kr).
\label{eq:psi_B_general}
\end{equation}
Then Eq.~\eqref{eq:inserted_general} becomes
\begin{equation}
S_\ell^{AB}(k_1,k_2)
= \left(\frac{2}{\pi}\right)^2\mathcal N_\ell^{AB}
\int_0^{\infty} dq\,q^2\,P_{AB}(q)
\Psi_{\ell}^{A}(k_1,q)\,\Psi_{\ell}^{B}(k_2,q).
\label{eq:sfb_general_psi_form}
\end{equation}

What this step does:
\begin{itemize}
\item $\Psi_\ell^X$ maps one latent mode $q$ into one observed SFB mode $k$ for
field $X$.
\item Physically, $\Psi_\ell^X$ is the response of the survey-weighted observable
to radial mode $q$.
\item Numerically, this isolates a reusable kernel that can be tabulated.
\end{itemize}

\subsection{Step 5: Window-function form}

If one defines convention-dependent windows $\mathcal W_\ell^X(k,q)$ from
$\Psi_\ell^X(k,q)$ (for example by absorbing factors of $2/\pi$, $q$, and $k$),
the same spectrum is written as
\begin{equation}
S_\ell^{AB}(k_1,k_2)
= \int_0^{\infty} dq\,P_{AB}(q)
\mathcal W_\ell^A(k_1,q)\,\mathcal W_\ell^B(k_2,q)
\times \mathcal F_\ell^{AB}(k_1,k_2,q),
\label{eq:sfb_window_general}
\end{equation}
where $\mathcal F_\ell^{AB}$ tracks whichever prefactors are not absorbed into
the window definition. In the common SuperFaB convention, all such prefactors
are absorbed so that $\mathcal F_\ell^{AB}=1$.

What this step does:
\begin{itemize}
\item It gives the compact estimator-friendly form used in SFB pipelines.
\item It makes normalization choices explicit and easy to compare between codes.
\item It separates physics content from convention choices.
\end{itemize}

\subsection{Why the single-Bessel transform is not sufficient}

A single transform,
\begin{equation}
\widetilde W_\ell^X(k)=4\pi\int_0^{\infty}dr\,r^2W^X(r)j_\ell(kr),
\end{equation}
is useful but does not yet contain the coupling between observed mode $k$ and
latent mode $q$. The full SFB covariance needs the two-Bessel object
$\Psi_\ell^X(k,q)$, because each observed mode receives contributions from a
continuum of latent $q$ modes once selection effects and projection are present.

\subsection{Specialization I: Galaxy clustering ($gg$)}

For galaxy clustering:
\begin{equation}
\alpha_g=0,
\qquad
\mathcal N_\ell^{gg}=\frac{2}{\pi},
\qquad
W^g(r)=b(r)n(r)D(r).
\end{equation}
Hence
\begin{equation}
C_\ell^{gg}(r_1,r_2)
=\frac{2}{\pi}W^g(r_1)W^g(r_2)
\int_0^{\infty}dq\,q^2P_{gg}(q)j_\ell(qr_1)j_\ell(qr_2),
\end{equation}
and
\begin{equation}
\Psi_\ell^g(k,q)=\int_0^{\infty}dr\,r^2W^g(r)j_\ell(kr)j_\ell(qr).
\end{equation}
Therefore
\begin{equation}
S_\ell^{gg}(k_1,k_2)
=\left(\frac{2}{\pi}\right)^3
\int_0^{\infty}dq\,q^2P_{gg}(q)
\Psi_\ell^g(k_1,q)\Psi_\ell^g(k_2,q).
\label{eq:sfb_gg_final}
\end{equation}

Interpretation for $gg$:
\begin{itemize}
\item $b(r)$ weights matter fluctuations into galaxy fluctuations.
\item $n(r)$ encodes the survey radial distribution.
\item $D(r)$ adds growth evolution.
\item Off-diagonal $k_1\neq k_2$ terms quantify radial mode coupling induced by
finite radial support and evolution.
\end{itemize}

\subsection{Specialization II: Shear ($ss$)}

For shear in the beyond-Limber kernel form:
\begin{equation}
\alpha_s=2,
\qquad
\mathcal N_\ell^{ss}=\frac{2(\ell+2)!}{\pi(\ell-2)!},
\end{equation}
with $W^s(r)$ the lensing efficiency kernel (including geometric and background
factors, such as $a^{-1}(r)$ and line-of-sight source weighting).

Then
\begin{equation}
C_\ell^{ss}(r_1,r_2)
=\mathcal N_\ell^{ss}W^s(r_1)W^s(r_2)
\int_0^{\infty}dq\,q^2P_{mm}(q)
\frac{j_\ell(qr_1)}{(qr_1)^2}
\frac{j_\ell(qr_2)}{(qr_2)^2},
\end{equation}
and
\begin{equation}
\Psi_\ell^s(k,q)=\int_0^{\infty}dr\,r^2W^s(r)
\frac{j_\ell(qr)}{(qr)^2}j_\ell(kr).
\end{equation}
Therefore
\begin{equation}
S_\ell^{ss}(k_1,k_2)
=\left(\frac{2}{\pi}\right)^2\mathcal N_\ell^{ss}
\int_0^{\infty}dq\,q^2P_{mm}(q)
\Psi_\ell^s(k_1,q)\Psi_\ell^s(k_2,q).
\label{eq:sfb_ss_final}
\end{equation}

Interpretation for $ss$:
\begin{itemize}
\item The $(qr)^{-2}$ factors encode the spin-2 and derivative structure of
lensing observables in this kernel representation.
\item The factorial prefactor boosts high-$\ell$ behavior consistently with
spin-2 harmonic normalization.
\end{itemize}

\subsection{Specialization III: Cross-correlation ($gs$ or $sg$)}

For galaxy-shear cross correlation:
\begin{equation}
\alpha_g=0,
\qquad
\alpha_s=2,
\end{equation}
with a normalization $\mathcal N_\ell^{gs}$ fixed by the exact convention used
in the chosen theory pipeline.

The radial kernel is
\begin{equation}
C_\ell^{gs}(r_1,r_2)
=\mathcal N_\ell^{gs}W^g(r_1)W^s(r_2)
\int_0^{\infty}dq\,q^2P_{gm}(q)
j_\ell(qr_1)\frac{j_\ell(qr_2)}{(qr_2)^2},
\end{equation}
which yields
\begin{equation}
\Psi_\ell^g(k,q)=\int_0^{\infty}dr\,r^2W^g(r)j_\ell(kr)j_\ell(qr),
\end{equation}
\begin{equation}
\Psi_\ell^s(k,q)=\int_0^{\infty}dr\,r^2W^s(r)\frac{j_\ell(qr)}{(qr)^2}j_\ell(kr),
\end{equation}
and
\begin{equation}
S_\ell^{gs}(k_1,k_2)
=\left(\frac{2}{\pi}\right)^2\mathcal N_\ell^{gs}
\int_0^{\infty}dq\,q^2P_{gm}(q)
\Psi_\ell^g(k_1,q)\Psi_\ell^s(k_2,q).
\label{eq:sfb_gs_final}
\end{equation}

Interpretation for $gs$:
\begin{itemize}
\item One leg probes biased galaxy density, the other probes projected tidal
distortions from matter.
\item Asymmetry between the two legs appears directly through
$\alpha_g\neq\alpha_s$ and $W^g\neq W^s$.
\item Swapping the probe order gives $S_\ell^{sg}(k_1,k_2)$ with the two response
kernels exchanged.
\end{itemize}

\subsection{Practical summary for implementation}

To compute any channel $AB$ in practice:
\begin{enumerate}
\item Choose $(\alpha_A,\alpha_B,\mathcal N_\ell^{AB})$ for that observable pair.
\item Build $W^A(r)$ and $W^B(r)$ on the survey radial grid.
\item Tabulate $\Psi_\ell^A(k,q)$ and $\Psi_\ell^B(k,q)$.
\item Integrate Eq.~\eqref{eq:sfb_general_psi_form} over $q$ to obtain
$S_\ell^{AB}(k_1,k_2)$.
\item Optionally repackage into the compact window form
Eq.~\eqref{eq:sfb_window_general}.
\end{enumerate}

This closes the general derivation and its direct specialization to
$(gg)$, $(ss)$, and $(gs)$.
